%\documentclass{beamer}
%Для защит онлайн лучше использовать разрешение 16x9
\documentclass[aspectratio=169]{beamer}

% !TeX spellcheck = ru_RU
% !TEX root = vkr.tex
% Опциональные добавления используемых пакетов. Вполне может быть, что они вам не понадобятся, но в шаблоне приведены примеры их использования.
\usepackage{tikz} % Мощный пакет для создание рисунков, однако может очень сильно замедлять компиляцию
\usetikzlibrary{decorations.pathreplacing,calc,shapes,positioning,tikzmark}

% Библиотека для TikZ, которая генерирует отдельные файлы для каждого рисунка
% Позволяет ускорить компиляцию, однако имеет свои ограничения
% Например, ломает пример выделения кода в листинге из шаблона
% \usetikzlibrary{external}
% \tikzexternalize[prefix=figures/]

\newcounter{tmkcount}

\tikzset{
  use tikzmark/.style={
    remember picture,
    overlay,
    execute at end picture={
      \stepcounter{tmkcount}
    },
  },
  tikzmark suffix={-\thetmkcount}
}

\usepackage{booktabs} % Пакет для верстки "более книжных" таблиц, вполне годится для оформления результатов
% В шаблоне есть команда \multirowcell, которой нужен этот пакет.
\usepackage{multirow}
\usepackage{siunitx} % для таблиц с единицами измерений

\newcommand{\cd}[1]{\texttt{#1}}
\newcommand{\inbr}[1]{\left<#1\right>}

% Для названий стоит использовать \textsc{}
\newcommand{\OCaml}{\textsc{OCaml}}
\newcommand{\miniKanren}{\textsc{miniKanren}}
\newcommand{\BibTeX}{\textsc{BibTeX}}
\newcommand{\vsharp}{\textsc{V$\sharp$}}
\newcommand{\fsharp}{\textsc{F$\sharp$}}
\newcommand{\csharp}{\textsc{C$\sharp$}}
\newcommand{\GitHub}{\textsc{GitHub}}
\newcommand{\SMT}{\textsc{SMT}}

\newcolumntype{L}[1]{>{\raggedright\let\newline\\\arraybackslash\hspace{0pt}}m{#1}}
%\newcolumntype{C}[1]{>{\centering\let\newline\\\arraybackslash\hspace{0pt}}m{#1}}
\newcolumntype{R}[1]{>{\raggedleft\let\newline\\\arraybackslash\hspace{0pt}}m{#1}}

\usepackage{xcolor}

%  Команды и пакеты, не используемые в шаблоне, которые тем не менее могут быть полезными.

% \newcolumntype{Y}{>{\centering\arraybackslash}X}

% \usepackage{mathrsfs}

\lstdefinestyle{fsharp}{
  basicstyle=\small\ttfamily,
  keywordstyle=\color{blue},
  commentstyle=\color{lightblue},
  stringstyle=\color{red},
  numbers=left,
  numberstyle=\tiny\color{gray},
  stepnumber=1,
  numbersep=10pt,
  backgroundcolor=\color{white},
  showspaces=false,
  showstringspaces=false,
  showtabs=false,
  tabsize=2,
  breaklines=true,
  breakatwhitespace=false,
  morekeywords={async, let, rec, if, then, else, do, yield, return, use, match, with, function}
  literate={->}{{$\to$}}3 {==}{{$\equiv$}}1 {=/=}{{$\not\equiv$}}1 {|>}{{$\triangleright$}}3 {\\/}{{$\vee$}}2 {/\\}{{$\wedge$}}2 {>=}{{$\ge$}}1 {<=}{{$\le$}}
}


% То, что в квадратных скобках, отображается внизу по центру каждого слайда. 
\title[GPU Image Processing]{Реализация библиотеки по обработке изображений с использованием GPU для платформы .NET}

% То, что в квадратных скобках, отображается в левом нижнем углу. 
\institute[СПбГУ]{}

% То, что в квадратных скобках, отображается в левом нижнем углу.
\author[Савельева Полина]{Савельева Полина Андреевна, группа 22.Б07-мм}
 
\begin{document}
{
\setbeamertemplate{footline}{}
% Лого университета или организации, отображается в шапке титульного листа
\begin{frame}
  \includegraphics[width=1.4cm]{pictures/SPbGU_Logo.png}
\vspace{-35pt}
\hspace{-10pt}
\begin{center}
   \begin{tabular}{c}
        \scriptsize{Санкт-Петербургский государственный университет} \\
        %\scriptsize{Кафедра системного программирования}
    \end{tabular}
\titlepage
\end{center}

\btVFill

{\scriptsize
  % У научного руководителя должна быть указана научная степень
   \textbf{Научный руководитель:} к.~ф.-м.~н. С.~В.~Григорьев, доцент кафедры информатики \\
  % Консультанта может и не быть. Должна быть указана должность или ученая степень
   %\textbf{Консультант:}  П.П. Петров, программист ЗАО \enquote{Компания с ну очень-очень-очень длинным названием}\\
  % Для курсовой не обязателен. Должна быть указана должность или ученая степень
   %\textbf{Рецензент:} д.т.н., проф. И.И. Иванов, исполнительный директор ООО \enquote{Рога и копыта}  
 }
\begin{center}
  \vspace{5pt}
  \scriptsize{Санкт-Петербург\\
                 2023}
  \end{center}

\end{frame}
}

\begin{frame}[fragile]  
  \frametitle{Введение}
  \begin{itemize}
    \item Обработка изображений используется в различных сферах, таких как компьютерное зрение, медицина, графика и дизайн
    \item Использование GPU может значительно ускорить обработку изображений
    \item Популярная платформа .NET предоставляет инструменты для высокоуровневого программирования, в том числе для обработки изображений 
    \item Идея --- реализовать средство, позволяющее эффективно использовать возможности GPU при работе с изображениями на платформе .NET
  \end{itemize}
\end{frame}
            
\begin{frame}  
  \frametitle{Существующие решения}
    \begin{itemize}
    
        \item \textbf{Magick.NET} является .NET--оберткой для \textbf{ImageMagick}, позволяющая разработчикам использовать функциональность ImageMagick в приложениях, без необходимости вызова командной строки. Она поддерживает множество операций над изображениями, но только на CPU (в отличии от ImageMagick).
    
        \item \textbf{ImageSharp}: библиотека помимо основных инструментов по работе с изображениями предоставляет своего рода строительные блоки, с помощью которых пользователи могут разрабатывать дополнительные функции. Не имеет встроенных функций для работы с GPU.

        \item .NET Библиотеки для обработки на GPU в большинстве являются обёртками, которые сложно сопровождать, или ориентированы на отдельных производителей видеокарт
    \end{itemize}
    
\end{frame}

% Обязательный слайд: четкая формулировка цели данной работы и постановка задачи
% Описание выносимых на защиту результатов, процесса или особенностей их достижения и т.д.
\begin{frame}
  \frametitle{Постановка задачи}
  \textbf{Целью} работы является реализация библиотеки по обработке изображений с использованием GPU на платформе .NET %озвученной выше 
  
  \textbf{Задачи}:
  \begin{itemize}
    \item Реализовать следующие возможности:
    \begin{itemize}
        \item Сохранение в различных форматах
        \item Изменение размера (reseize)
        \item Обрезка (crop)
        \item Добавление водяного знака (watermark)
    \end{itemize}
    \item Сравнить производительности текущего решения и аналогов 
    \item Оформить и опубликовать nuget пакет
  \end{itemize}
\end{frame}

\begin{frame}
  \frametitle{Brahma.FSharp}
    \begin{itemize}
        \item Для сообщения с GPU и решения поставленных задач  необходим .NET инструмент, поддерживающий широкий список графических процессоров и позволяющий писать код только на высокоуровневых языках
        \item  Brahma.FSharp --- F\#~инструмент для работы с GPU различных производителей без использования низкоуровневых языков
    \end{itemize}
\end{frame}

\end{document}
